% Requires amsmath, amssymb, graphicx. Compile as part of a larger document;
% figure paths are relative to crates/tt-fingerprint-optimizer.
\doclearpage
\section{Pre-Filtering Optimization}
\label{sec:fp-opt}

A \textsc{like} filter over a string column is preceded by a fingerprint
pre-filter. Each row commits a bit vector of the features it contains,
and a row survives only if its vector contains the bits the prover
chose to test from the pattern's. The survivors are compacted before the
expensive \textsc{like} proof runs. The work splits between two parties.
At commit time the data owner assigns the column's features to bins,
from the column alone, and commits every row's fingerprint. At proving
time the prover, which holds the data, chooses per query which of the
pattern's bins to test, pricing each choice with a cost model of the
proof. On the benchmark column the pre-filter makes selective queries
$10.4\times$ faster to prove.

\subsubsection{Fingerprints}

\paragraph{Columns and rules.}
Every string column of every table is fingerprinted separately; columns
of other types are never fingerprinted. A column $\gamma$ is a corpus of
rows $R_\gamma$, each row a byte string, with $m_\gamma$ characters in
total. Its \emph{rule}, the assignment of features to bins defined
below, is stored as trusted table metadata in the table's
\texttt{.oracle}. The verifier reads the rule from the oracle, never from
the proof, so an untrusted prover cannot influence how a column is
fingerprinted. A data owner who declares one rule and commits limbs built
under another could hide rows of its own table, but it could equally
commit different data, so the rule grants it no new power.

\paragraph{Features.}
\label{par:fp-features}
For a byte string $r$ the feature set $\mathrm{feat}(r)$ contains every
byte (\emph{char}); every length-$2$ and length-$3$ window of $r$ read
over a $64$-symbol alphabet --- the $62$ alphanumerics, the space, and
one symbol for every other byte (\emph{bigram}, \emph{trigram}); and
every ordered pair $(a, b)$ of alphanumeric bytes such that the first
occurrence of $a$ lies before the last occurrence of $b$
(\emph{precedence}). Reading $n$-grams across spaces and punctuation lets
a pattern say where a word starts or ends: \texttt{\%l\%n\ \%} has the
bigram \texttt{"n\ "}, which only rows with a word ending in \texttt{n}
contain (\S\ref{par:fp-boundary}). A pattern $p = \%f_1\%\cdots\%f_n\%$
has the chars and $n$-grams of each literal factor on its own --- a
wildcard can match anything, so $n$-grams never cross factors --- plus
the precedence pairs of the concatenation $f_1 \cdots f_n$. Whenever $r$
matches $p$ the factors occur in $r$ in order and without overlap, so
\[
  \mathrm{feat}(p) \subseteq \mathrm{feat}(r).
\]
On the benchmark column \texttt{lineitem.l\_comment} ($|R| = 299{,}814$
rows, $m = 7{,}949{,}138$ characters) $2{,}203$ distinct features occur
in some but not every row: $35$ chars, $302$ bigrams, $1{,}138$ trigrams
and $728$ precedence pairs.

\paragraph{Fingerprints and bins.}
A rule is a width $B_\pi \le 2{,}048$ and a partial function
$\pi : \mathcal{F} \rightharpoonup \{0, \dots, B_\pi - 1\}$ from features
to bins; unassigned features set no bin. The row fingerprint
$\mathrm{fp}_\pi(r) \in \{0,1\}^{B_\pi}$ has bit $b$ set iff $r$ has some
feature $f$ with $\pi(f) = b$, so a bin is the \textsc{or} of its
features. The pattern's bins are $\varphi(p) = \pi(\mathrm{feat}(p))$. A
pre-filter tests a subset $Q \subseteq \varphi(p)$ and keeps the rows
with $Q \subseteq \mathrm{fp}_\pi(r)$. By the containment above this is
\emph{sound for every $\pi$ and every $Q$}: sharing bins and leaving bins
untested only add false positives, never drop a match, so both choices
are purely performance questions. Writing $M_b = \bigcup_{f : \pi(f) = b}
\mathrm{rows}(f)$ for the rows setting bin $b$, the survivors are
\[
  S_\pi(p, Q) \;=\; \bigcap_{b \in Q} M_b .
\]

\paragraph{Commitment and the pre-filter check.}
The data owner commits each fingerprint once, as $\lceil B_\pi / 8
\rceil$ separate columns of $8$-bit limbs; limb $j$ holds bins $8j$ to
$8j + 7$. A pre-filter testing $Q$ works only on the \emph{needed} limbs
$N = \{ \lfloor b/8 \rfloor : b \in Q \}$: one batched
bitwise-\textsc{and} lookup per needed limb into one shared $2^{16}$-row
table (table membership doubles as the range check), plus zero checks
that every kept row passes the subset test and a non-zero check that
every dropped row fails it. The comparison folds the needed limbs with a
challenge $\rho$ drawn after the masked limbs are committed,
$\sum_i \rho^i g_{N_i}$ against the same fold of $Q$. The limbs are
bytes, so two differing limb vectors agree under the fold with
probability at most $(|N| - 1)/|\mathbb{F}|$, for any number of limbs,
and the width is not bounded by the scalar modulus.

A limb no tested bin lies in is never read. We measured this directly:
\texttt{\%erve\%} touches $8$ limbs whether $32$ or $64$ are committed,
and proves in $61.4$ and $59.4$\,s respectively (median of three).
Committing wider costs the data owner --- the oracle grows from $34.2$ to
$62.4$\,KB, and the extra columns must be encoded and committed --- but
costs the prover nothing.

\subsubsection{Rules at Commit Time}
\label{sec:fp-modes}

The prover chooses its bins per query (\S\ref{par:fp-select}), so the
data owner's rule decides only which features get a bin and which share
one. A shared bin costs selectivity, since a row passes if it has
\emph{any} of the bin's features. No rule looks at queries: each is built
from the column's active rows alone. The data owner picks one of two
modes for a table (\texttt{commit --fp-mode}, or the
\texttt{TT\_FP\_MODE} environment variable):
\begin{itemize}
  \item \emph{cluster} (the default): every feature, grouped into $512$
        bins by the rows the features share (\S\ref{sec:fp-cluster});
  \item \emph{single}: one private bin per feature, the most common
        first, up to $2{,}048$ (\S\ref{sec:fp-single}).
\end{itemize}
Every string column of the table gets its own rule under the chosen
mode. In both modes a feature present in every row gets no bin, since it
can never remove a row, and a column of fewer than $2^{14}$ rows gets a
rule of width zero and no pre-filter, since its padded domains are too
small for a pre-filter to repay its limbs.

\paragraph{Cluster rules.}
\label{sec:fp-cluster}
A query for a feature $f$ in bin $B$ keeps every row of $\bigcup B$
where only the rows of $f$ could match, so sharing a bin wastes the rows
that hold the bin's other features but not $f$. Weighting each feature by
how often a query contains it, a bin wastes
\[
  W(B) \;=\; \sum_{f \in B} w_f \Bigl(\bigl|\textstyle\bigcup B\bigr| - |\mathrm{rows}(f)|\Bigr)
\]
rows. With no queries to count, $w_f = |\mathrm{rows}(f)| / |R|$: a
pattern is cut from the column's own text, so it contains a feature about
as often as the rows do. Merging bins $A$ and $B$ then adds
\[
  \Delta(A, B) \;=\; (w_A + w_B)\,|A \cup B| - w_A |A| - w_B |B|
\]
wasted rows, with $w_A$ a bin's summed weight and $|A|$ the rows of its
union. $\Delta$ is small in exactly two cases: features that occur
together (\texttt{th}, \texttt{the} and \texttt{he}), whose union adds
almost no row to either, and rare features, whose union is small. It is
large for a common feature and a rare one, which would hand the rare
feature's queries all the common feature's rows.

The rule starts from one bin per feature ($2{,}203$ on
\texttt{l\_comment}; at most the $4{,}096$ most common on a column with
more) and repeatedly makes the cheapest merge until $512$ bins remain. A
table of all pairwise $\Delta$ is filled once, in parallel, and after a
merge only the new bin's row is recomputed; each bin keeps its cheapest
partner, so a merge costs one pass over the bins. Every feature keeps a
bin. Columns longer than $2^{20}$ rows are sampled at an even stride.
Width is the only parameter: along the merge path $J$
(\S\ref{sec:fp-results}) falls quickly and then flattens, and $512$ bins,
$64$ limbs per row, is $0.5$\,s of $J$ from the widest rule at a quarter
of its limbs.

\paragraph{Single-feature rules.}
\label{sec:fp-single}
A private bin per feature leaves the prover every choice: no feature's
rows are blurred by another's. The only choice left is which features
get a bin when a column has more than the $2{,}048$ bins a fingerprint
may use, as all three large comment columns do. The rule keeps those in
the most rows, for the reason behind the cluster rule's weights. Keeping
the rarest instead is clearly worse, and keeping the most common matches
a rule that gives a bin to exactly the features the workload uses
(modeled $J$ under the earlier fit of the cost model, bins chosen per
query):
\begin{center}
\begin{tabular}{lrrr}
  column & most common first & rarest first & features the workload uses \\
  \hline
  \texttt{l\_comment}  & 205.7 & 227.1 & 205.7 \\
  \texttt{o\_comment}  & 123.2 & 124.6 & 123.2 \\
  \texttt{ps\_comment} & 165.2 & 166.2 & 165.2
\end{tabular}
\end{center}
The rule is as wide as its features allow, $2{,}048$ bins or $256$
committed limbs per row on these columns.

\paragraph{Why not maximize entropy.}
\label{sec:fp-entropy}
A natural workload-free alternative picks the bins that carry the most
information about the rows. Built in order, each bin is the
\textsc{or} of the features that greedily maximize its conditional
entropy given the bins before it,
$H(Y_b \mid Y_1, \dots, Y_{b-1}) = \sum_C \frac{|C|}{|R|}\,
h_2\bigl(|C \cap U_b| / |C|\bigr)$ over the cells $C$ of rows sharing a
fingerprint so far --- decision-tree induction with \textsc{or}-merged
splits. We implemented this rule as a third commit mode and removed it;
the numbers below are modeled under the earlier fit of the cost model,
where the default rule had $J = 206.7$\,s. It works: it is as sound as
any rule, and it lowers $J$ from $327.7$ to $249.7$\,s. But it optimizes the wrong quantity. Entropy rewards telling
rows apart, while the prover pays for the rows a pattern keeps. A bin is
most informative when it splits the rows near half and half, so the rule
favors features present in about half the rows, which separate rows well
but prune a selective pattern's candidates poorly; the rare features a
selective pattern needs carry little entropy and get no bin. It also
stops growing once the rows are told apart: on \texttt{l\_comment} after
$122$ bins, whatever width is asked for, with the same $122$ when
trigrams are offered too. Below $128$ bins its shared bins let it lead
the single-feature rule ($295.2$, $275.9$ and $263.3$\,s at $16$, $32$
and $64$ bins, against $305.3$, $290.9$ and $271.5$), but the cluster
rule beats it at every width ($275.6$, $246.8$ and $234.2$\,s), and at
its full width it trails the default by $43$\,s of $J$, nearly all of it
in the selective regimes ($95.7$ against $44.0$\,s per query below $1\%$,
$179.4$ against $88.0$ between $1$ and $20\%$). It is also the slowest to
build: $7$\,s on \texttt{l\_comment}, against $3$\,s for the cluster
rule.

\paragraph{Packing bins into limbs.}
Both rules number their bins by rows, most first, and limb $j$ takes
bins $8j$ to $8j + 7$; which features a pattern holds together plays no
part. The order does not matter: a random order prices the same to
within $0.01$\,s of $J$ at every width tested. Nor would a better one pay.
Packing together the bins that the evaluation workload's own queries
want tested --- an in-sample layout no commit-time rule could build ---
cuts the limbs a selective query touches from $3.8$ to $2.9$ on the
default rule, yet gains only $0.15$\,s of $J$ under the earlier fit of
the cost model ($206.67 \to 206.52$\,s;
$0.12$--$0.18$\,s for the cluster rule at $128$--$1{,}024$ bins and the
single-feature rule at $512$ and $2{,}048$). By the time a selective
query has tested its first few limbs its survivors are near the matches,
so a saved limb is worth one $c_{\mathrm{limb}}$ in one regime of six.

\paragraph{Commit cost.}
Committing the $2^{19}$-row \texttt{lineitem} benchmark table (building
every string column's rule, encoding and committing, not loading the
key; medians of three) takes $1.15$\,s with no fingerprints, $4.3$ to
$4.9$\,s with cluster rules of $16$ to $2{,}048$ bins, and $1.3$ to
$2.5$\,s with single-feature rules of the same widths
(Figure~\ref{fig:fp-widths}, right). Building the cluster rule is most of
it, about $3$\,s at any width, since the merge path runs from every
feature down to the width asked; building the single-feature rule takes
$0.2$\,s. The limb commitments themselves are cheap, about $5$\,ms per
committed limb column. The oracle grows from $6.2$\,KB to $54$\,KB at the
default width and to $87$--$89$\,KB at $2{,}048$ bins, mostly the rules
themselves.

\subsubsection{Choosing Bins at Proving Time}
\label{par:fp-select}

The prover chooses $Q$ per query, in whole limbs, because the limb is
what the pre-filter pays for: a touched limb costs one lookup whether it
tests one of the pattern's bins or all of them. At plan time it
fingerprints the column's active rows and counts the pattern's matches
$t$. Starting from every row, it repeatedly adds the limb, among those
holding bins of $\varphi(p)$, whose test (every bin of $\varphi(p)$ in
that limb) leaves the fewest survivors, ties to the lowest limb, and
after each step prices the query with the cost model of
\S\ref{par:fp-cost}. It keeps the cheapest prefix of that chain, ties to
the shorter, possibly the empty one, in which case the pre-filter is
removed from the plan. Each step is one bitset intersection per
remaining limb, and a pattern touches at most a few dozen limbs, so the
choice takes milliseconds.

The whole chain is priced rather than stopping at the first limb that
does not pay. The cost falls in steps at powers of two
(\S\ref{par:fp-cost}), so a limb that saves nothing can come just before
one that crosses a step: on \texttt{l\_comment} that happens for $99$ of
the $1{,}239$ workload patterns, and stopping at the first rise would
cost $3.1$\,s per query on average.

The chosen list travels with the proof as a plan hint. The verifier
checks that it is a non-empty subset of $\varphi(p)$, which it computes
from the public pattern and the rule in the oracle, and builds the
pre-filter from it. Nothing else about the choice needs to be agreed:
any subset is sound, and the cost model is a heuristic whose errors cost
proving time, never correctness.

Testing a limb pays only while it removes rows, so the speed-up against
the number of limbs tested rises and then falls
(Figure~\ref{fig:fp-chain}): it climbs steeply over the first few limbs,
then drops as each extra limb adds about $c_{\mathrm{limb}}$ and removes
almost nothing. On the default rule the prover stops at $3.1$ limbs on
average for patterns matching under $1\%$ of rows, $1.8$ between $1$ and
$20\%$, and $0.2$ between $40$ and $60\%$; above $60\%$ it never tests a
limb.

\begin{figure}
  \centering
  \includegraphics[width=\linewidth]{results/chain-cluster-l_comment.png}
  \caption{The prover's limb choice: speed-up over no pre-filter after
    the first $k$ limbs of each query's greedy chain, per selectivity
    regime (mean no-filter cost over mean cost), on the default rule
    ($512$ bins) over the \texttt{l\_comment} workload, modeled. A query
    with fewer candidate limbs keeps its cost at all of them. Each panel's
    axis is zoomed on the top of its curve.}
  \label{fig:fp-chain}
\end{figure}

\subsubsection{Prover Cost Model}
\label{par:fp-cost}

\paragraph{The model.}
A query's proving time is modeled from its \emph{shape}: input rows
$|R|$ and characters $m$, survivors $s$ and their characters $m_s$,
matches $t$, touched limbs $\lambda$, total literal length $L$ and
factor count $n$. The shape fixes what the planner does. Write $D(x)$ for
the next power of two at or above $x$ and $D^+(x)$ for the strict one
(doubling an exact power). Tables are padded to a power of two and the
planner counts the padding, so a filter over a scan is compared against
$T = D(|R|)$. Then
\begin{align*}
  \mathrm{remat}  &= [\lambda > 0] \wedge [D^+(s) < D^+(T)], &
  (R', m', P)     &= \begin{cases} (s, m_s, s) & \mathrm{remat},\\
                                    (|R|, m, T) & \text{otherwise,} \end{cases}\\
  \mathrm{second} &= [D^+(t) < D^+(P)], &
  m_t             &= t\, m' / R'.
\end{align*}
The survivors are compacted (``rematerialized'') iff remat, and the
\textsc{like} runs on $(R', m')$. Every compaction of a string column
carries a data-preserving update check (DPUC) binding the compacted
characters to their source rows, so a compaction costs work over its
\emph{character} domain; the DPUC reads characters only, never
fingerprint limbs. The model is linear in fifteen regressors, each priced
on the padded domain the stage actually runs on:
\begin{align*}
  C = {}& a_0
    && \text{(fixed floor of every proof)} \\
  &+ a_{\mathrm{pf}}\, [\lambda > 0]
    && \text{(pre-filter check, paid once if any bin is tested)} \\
  &+ c_{\mathrm{limb}}\, \lambda
    && \text{(AND lookup per touched limb)} \\
  &+ \mathrm{remat} \cdot a_r
    && \text{(fixed cost of compacting the survivors)} \\
  &+ \mathrm{remat} \cdot a_{rs}\, D(s)
    && \text{(compaction, per surviving row)} \\
  &+ \mathrm{remat} \cdot a_{rc}\, D(m_s)
    && \text{(compaction DPUC, per surviving character)} \\
  &+ a_c\, D(m')
    && \text{(\textsc{like}, per input character)} \\
  &+ a_{cL}\, (L - 1)\, D(m')
    && \text{(\textsc{like}, per character per extra literal character)} \\
  &+ a_{cn}\, (n - 1)\, D(m')
    && \text{(\textsc{like}, per character per extra factor)} \\
  &+ a_w\, D(R')
    && \text{(\textsc{like}, per input row)} \\
  &+ a_{wL}\, (L - 1)\, D(R')
    && \text{(\textsc{like}, per row per extra literal character)} \\
  &+ a_{wn}\, (n - 1)\, D(R')
    && \text{(\textsc{like}, per row per extra factor)} \\
  &+ \mathrm{second} \cdot \bigl(a_2 + a_{2r}\, D(R') + a_{2c}\, D(m_t)\bigr)
    && \text{(see below).}
\end{align*}
Because every term reads a padded domain, $C$ is a \emph{staircase}:
removing survivors buys nothing until the surviving domain crosses a
power of two, then buys a whole step.

\begin{center}
\begin{tabular}{llr@{\qquad}llr}
  term & & fitted & term & & fitted \\
  \hline
  floor & $a_0$ & $17.88$ s &
  \textsc{like} chars & $a_c$ & $4.21 \times 10^{-6}$ \\
  pre-filter & $a_{\mathrm{pf}}$ & $0$ &
  \quad per extra literal char & $a_{cL}$ & $3.07 \times 10^{-6}$ \\
  touched limb & $c_{\mathrm{limb}}$ & $0.659$ s$^\dagger$ &
  \quad per extra factor & $a_{cn}$ & $8.42 \times 10^{-6}$ \\
  remat & $a_r$ & $10.26$ s &
  \textsc{like} rows & $a_w$ & $6.21 \times 10^{-5}$ \\
  \quad rows & $a_{rs}$ & $1.05 \times 10^{-4}$ &
  \quad per extra literal char & $a_{wL}$ & $0$ \\
  \quad chars & $a_{rc}$ & $6.65 \times 10^{-6}$ &
  \quad per extra factor & $a_{wn}$ & $0$ \\
  second & $a_2$ & $0$ &
  second rows & $a_{2r}$ & $0$ \\
  & & & second chars & $a_{2c}$ & $1.51 \times 10^{-6}$
\end{tabular}
\end{center}
Domain terms are in seconds per padded entry. $^\dagger$Held fixed; see
below. The three ``second'' terms are named for a compaction of the
\textsc{like} output that the planner does not perform (it leaves alone
any node that only projections separate from the result check); the fit
all but drops them, keeping only a small per-character term. The
per-stage split of the model is not a breakdown of the proof. The
coefficients are compiled into the prover
(\texttt{crates/tt-arithmetic/cost\_model.toml}), which uses them only to
choose bins.

\paragraph{Calibration.}
The coefficients other than $c_{\mathrm{limb}}$ are fitted by
non-negative least squares (\texttt{tt-fp-opt fit}) to $50$ release-mode
whole-prove timings on the benchmark column (\texttt{tt-fp-measure}),
taken on an otherwise idle machine and ranging from $28$ to $388$\,s. Two
workload queries per regime, drawn with a fixed seed among those with at
most two factors and at least one match, are proved with no
fingerprints; the six of them matching under $40\%$ of rows are proved
again against the table committed with the cluster rule at $16$, $64$,
$256$, $512$, $1{,}024$ and $2{,}048$ bins, each query's shape taken from
the prover's own choice. The fit has RMS error $3.2$\,s ($3.6$\,s
leave-one-out), a held-out error of $3.7\%$ on average and $11.2\%$ at
most. It replaces a fit to $69$ proofs taken on an earlier build of the
prover, which on these $50$ proofs predicts $1.2$--$1.6\times$ too long
(RMS $59$\,s) and overstates the speed-up of the $1$--$40\%$ regimes by
$25$--$35\%$. Comparisons marked ``under the earlier fit'' in this
section were modeled with it; they compare rules under one model, and
where we recomputed such comparisons under the refit (the width sweeps
below) their ranking did not change. The fitting set has at most two factors and ten
literal characters, so patterns with three or four factors extrapolate
the per-factor terms.

A free fit cannot recover $c_{\mathrm{limb}}$: a few seconds of limb cost
is lost against errors on proofs of several hundred seconds, and least
squares drives it to zero. It is therefore measured directly, by proofs
that hold the survivor set fixed and vary only the limbs. Moving whole
bins keeps which features \emph{share} a bin, so every row keeps its set
bits and the survivors are unchanged to the row, while the same bins are
spread over more limbs. Three such experiments, on the earlier build,
agree:
\begin{itemize}
  \item a limb series spreading one pattern's $65$ bins over $9$, $16$,
        $24$ and $32$ limbs ($11$ survivors throughout) gives $0.66$\,s
        per limb, all points within $0.2$\,s of the line;
  \item a rule restricted to $632$ features and relabelled
        $b \mapsto sb$ for $s = 1, \dots, 4$ gives \texttt{\%quickly final
        deposits wake\%} $16$, $31$, $42$ and $50$ touched limbs with the
        same $3$ survivors, proving in $44.3$, $53.3$, $59.8$ and
        $63.5$\,s: $0.570$\,s per limb, $R^2 = 0.998$;
  \item relabelling $b \mapsto 2b$ in another rule gives $0.60$ and
        $0.62$\,s per limb on two patterns ($2$ and $217$ survivors)
        between $29$ and $62$ limbs.
\end{itemize}
So $c_{\mathrm{limb}} \in [0.57, 0.66]$\,s and the price is linear in the
touched limbs up to at least $50$ of them. The model carries $0.659$, the
high end.

\paragraph{What the measurements show.}
Selective patterns gain most: \texttt{\%ely f\%nal d\%} ($377$ matches)
proves in $387.6$\,s with no pre-filter and $30.5$\,s against the default
rule ($12.7\times$), and \texttt{\%ach\%} ($1{,}494$ matches) in $138.5$
and $30.4$\,s ($4.6\times$). Less selective ones gain less, because the
compaction and the \textsc{like} over the survivors still cost most of a
proof: \texttt{\%refully\%} ($11\%$ of rows) $3.3\times$,
\texttt{\%ic\%} ($19\%$) $1.7\times$, \texttt{\%t\%y\%} and
\texttt{\%ly\%a\%} ($33$ and $36\%$) $1.4$ and $1.3\times$. A pre-filter
that removes too few rows would make the proof \emph{slower}, since it
pays for the check and the compaction without shrinking the
\textsc{like}; the prover avoids that by testing nothing, as it does for
every pattern above $60\%$ on the default rule. Past the point where
extra limbs remove rows, each one only adds $c_{\mathrm{limb}}$, which is
the fall on the right of Figure~\ref{fig:fp-chain}; the prover's chain
stops at the bottom of that curve per query.

\subsubsection{Results}
\label{sec:fp-results}

\paragraph{Workload and metric.}
Rules never see queries; the workload exists only to evaluate them. We
generate a workload $P_\gamma$ from the corpus itself, separately for
each of six selectivity regimes $\sigma \in [\theta_{i-1}, \theta_i)$
with $(\theta_i) = (0, 0.01, 0.2, 0.4, 0.6, 0.8, 1]$. Candidates are
$1$--$3$ ordered literal factors cut from a random row (so they match at
least that row), a share of short patterns (up to $4$ factors of
$1$--$2$ bytes, the only ones common enough for high selectivities), and
a share with scrambled or cross-row factors (so some match nothing).
Each candidate's selectivity is computed exactly, sampling is stratified
over equal-width bands within a regime, and a regime stops early when
the corpus runs out of distinct patterns for it. On
\texttt{l\_comment} this yields $300$, $300$, $300$, $265$, $62$ and $12$
patterns, using $1{,}238$ distinct features; the corpus has very few
distinct patterns matching more than $60\%$ of rows. We summarize a rule
by
\[
  J \;=\; \frac{1}{6} \sum_{i=1}^{6} \frac{1}{|P_i|} \sum_{p \in P_i} C(p),
\]
the modeled proving time $C(p)$ (\S\ref{par:fp-cost}) with the prover's
choice of bins, averaged per regime and then over regimes, so every
regime weighs the same. Offline pricing computes $|S|$ and its character
count exactly: each feature carries a row bitset, $M_b$ is a union of
bitsets, the intersection is a word-wise \textsc{and}, and the
characters are $\sum_j 2^j \lvert S \wedge P_j \rvert$ over the
bit-planes $P_j$ of the row lengths. The character count must be exact:
estimating it as survivors times the mean length understates it (a row
holding a pattern's features is longer than average, $29.7$ against
$26.5$ characters here), and the cost model reads it through a
power-of-two domain, so the estimate falls on the wrong side of a step
often enough to misprice a rule by $15$\,s of $J$ (under the earlier
fit). Without a pre-filter $J = 235.2$\,s on \texttt{l\_comment}.

\paragraph{The two modes.}
Each mode at the width it commits, bins chosen per query:
\begin{center}
\setlength{\tabcolsep}{5pt}
\begin{tabular}{lrrrr}
  selectivity & none & cluster, 512 & single, 2{,}048 & perfect \\
  \hline
  $< 1\%$       & 360.6 & \textbf{34.6}  &  33.1 &  30.5 \\
  $1$--$20\%$   & 207.8 & \textbf{74.9}  &  73.6 &  64.1 \\
  $20$--$40\%$  & 221.0 & \textbf{164.2} & 164.1 & 143.4 \\
  $40$--$60\%$  & 250.0 & \textbf{244.0} & 243.5 & 205.8 \\
  $60$--$80\%$  & 217.0 & \textbf{217.0} & 217.0 & 217.0 \\
  $80$--$100\%$ & 154.7 & \textbf{154.7} & 154.7 & 154.7 \\
  \hline
  $J$           & 235.2 & \textbf{148.2} & 147.7 & 135.9 \\
  limbs         &       & 64             & 256   &
\end{tabular}
\end{center}
Seconds per query, modeled; \emph{perfect} is a filter that keeps
exactly the matching rows at the cost of one limb, or no pre-filter when
that is cheaper, the ceiling for any fingerprint. The default (bold) is
$10.4\times$ faster than no pre-filter on patterns matching under $1\%$
of rows and $2.8\times$ between $1$ and $20\%$, $0.5$\,s of $J$ behind
the single-feature rule at a quarter of its committed limbs. Above $60\%$
no filter can help: even a perfect one pays for a compaction that saves
less than it costs.

\paragraph{Width by commit mode.}
Figure~\ref{fig:fp-sweep} walks the width axis once per mode, every
query choosing its own bins. The cluster curve takes the rule at each
point of its merge path, the single-feature curve keeps the $k$ most
common features. Modeled $J$ in seconds:
\begin{center}
\begin{tabular}{lrrrrrrrr}
  bins & 16 & 32 & 64 & 128 & 256 & 512 & 1024 & 2048 \\
  \hline
  cluster & 194.4 & 172.4 & 163.7 & 158.1 & 152.0 & 148.2 & 147.7 & 147.7 \\
  single  & 220.0 & 208.6 & 193.9 & 171.6 & 160.7 & 150.6 & 148.1 & 147.7 \\
\end{tabular}
\end{center}
The cluster rule is best at every width from $4$ bins to $1{,}024$; at
$2{,}048$ it and the single-feature rule coincide. Sharing bins is what
lets it cover every feature in few bins: at $128$ bins it is $13.5$\,s
ahead of the single-feature rule, which has only the $128$ most common
features, and it reaches $152.0$\,s with $256$ bins, a level the
single-feature rule reaches only between $384$ and $512$. Nearly all the
difference between modes is in the two selective regimes; above $40\%$
selectivity no rule gains more than $1.1\times$.

\begin{figure}
  \centering
  \includegraphics[width=\linewidth]{results/sweep-cluster-l_comment.png}\\[1ex]
  \includegraphics[width=\linewidth]{results/sweep-single-l_comment.png}
  \caption{Pre-filter speed-up per selectivity regime as the committed
    rule gets wider, for the cluster (top) and single-feature (bottom)
    commit modes. Every query tests the limbs the prover's chain picks.
    Speed-up over no pre-filter, modeled on the \texttt{l\_comment}
    workload; both axes are logarithmic and shared by the two panels.}
  \label{fig:fp-sweep}
\end{figure}

\paragraph{What width costs.}
Figure~\ref{fig:fp-widths} sets the prover's gain against what a wider
rule costs. The proof grows only when a query runs the pre-filter: the
fingerprint commitments are in the oracle, not the proof, and a query
that tests no limb proves exactly as without fingerprints (within
$40$\,bytes on every measured proof). Over the $50$ measured proofs the
size is, to a mean error of $3.8\%$ (leave-one-out, $14.7\%$ at worst),
\[
  14.7 + 4.5\,n + 0.8\,L + 7.8\,[\lambda > 0] \ \text{KB},
\]
a fixed $7.8$\,KB for the pre-filter check and the compaction when any
limb is tested; the number of limbs and the domain sizes add nothing to
the fit. Applied to the workload, the proof grows by at most $31\%$ in
any regime, by less in the regimes that rarely filter, and it stops
growing once the prover's choices settle, by $256$--$384$ bins for the cluster
rule. The data owner's cost is measured directly (\S\ref{sec:fp-modes}):
commit time is nearly flat for the cluster rule, whose merge dominates,
and grows about $2\times$ for the single-feature rule, while the oracle
grows to $8.5\times$ and $14\times$ at $512$ and $2{,}048$ bins, still
under $90$\,KB.

\begin{figure}
  \centering
  \includegraphics[width=\linewidth]{results/widths-cluster-l_comment.png}\\[1ex]
  \includegraphics[width=\linewidth]{results/widths-single-l_comment.png}
  \caption{What a wider committed rule buys and costs, for the cluster
    (top) and single-feature (bottom) rules on \texttt{l\_comment}.
    Left: prover speed-up per selectivity regime, modeled over the
    workload with the prover's own limb choice; circles are measured.
    Middle: proof size over the size with no pre-filter, from the fitted
    proof-size model. Right: measured commit time and oracle size over a
    commit with no fingerprints.}
  \label{fig:fp-widths}
\end{figure}

\paragraph{Word boundaries.}
\label{par:fp-boundary}
The largest single gain came from the features. With $n$-grams over
alphanumerics only, a space breaks every $n$-gram, so a pattern such as
\texttt{\%l\%n\ \%} (a word ending in \texttt{n}) or \texttt{\%\ i\%n\%}
(a word starting with \texttt{i}) falls back to single characters that
nearly every row has: \texttt{\%l\%n\ \%} then kept $179{,}530$ of
$299{,}814$ rows and took $283$\,s. Over the $64$-symbol alphabet the
same pattern has \texttt{"n\ "}, keeps $21{,}351$ rows, and proves in
$66$\,s. Across the workload the $<1\%$ regime gains $8$\,s per query and
the $1$--$20\%$ regime $16$\,s (under the earlier fit). Two further feature families were priced
and rejected: $4$-grams over the same alphabet add $0.7$\,s, and
character-count features (``at least two \texttt{e}'') add nothing,
because the characters involved are in nearly every row and the cost is
a staircase in the survivors.

\paragraph{What remains.}
The default rule is $4.1$\,s per query from a perfect filter in the
$<1\%$ regime, $10.8$\,s in the $1$--$20\%$ regime, and at the ceiling
in the two least selective, where no filter pays. Most of the gap sits in
the $20$--$60\%$ regimes ($20.8$ and $38.1$\,s), in two kinds of pattern
no fingerprint of this feature set can filter:
single-letter multi-factor patterns such as \texttt{\%e\%\ \%e\%s\%},
whose literals nearly every row contains, and patterns such as
\texttt{\%ca\%areful\%}, whose factors must not overlap while their
features cannot say so. Closing that gap needs a filter that understands
order and overlap, which is what the \textsc{like} proof itself does.

The workload is \texttt{workloads/l\_comment.csv} and the cost model
\texttt{crates/tt-arithmetic/cost\_model.toml}; \texttt{tt-fp-opt
sweep} reproduces every modeled number in the results from them, into
the tables in \texttt{results/}, and \texttt{scripts/plot\_sweep.py},
\texttt{plot\_chain.py} and \texttt{plot\_widths.py} the figures. The
measured proofs and commits are in \texttt{results/proofs-l\_comment.csv}
and \texttt{results/commits-l\_comment.csv} (from
\texttt{tt-fp-measure}), and the points the cost model is fitted to in
\texttt{results/cost-points-l\_comment.csv}.
